% Archivo principal: main.tex
\documentclass[12pt,letterpaper]{report}
\usepackage[justification=centering]{caption}

% Importar configuración externa
% Archivo de configuración: configuracion.tex

% Configuración de márgenes
\usepackage[left=3.5cm, right=3cm, top=3cm, bottom=3cm]{geometry}

% Configuración de idioma y tipografía
\usepackage[spanish]{babel}
\usepackage[T1]{fontenc}
\usepackage{times} % Times New Roman
\usepackage{fontsize}
\usepackage{xcolor}
\usepackage{geometry}
\usepackage{graphicx}
\usepackage{tocloft}

\graphicspath{{./img/}}
% Configuración del interlineado
\renewcommand{\baselinestretch}{1.5} % Interlineado 1.5

% Configuración de numeración de páginas
\usepackage{fancyhdr}
\pagestyle{fancy}
\fancyhf{}
\rfoot{\thepage} % Número de página alineado a la derecha
\renewcommand{\headrulewidth}{0pt} % Eliminar línea de encabezado

% Definición de macros para datos del documento
\newcommand{\titulo}[1]{\def\Titulo{#1}}
\newcommand{\autor}[1]{\def\Autor{#1}}
\newcommand{\tutor}[1]{\def\Tutor{#1}}
\newcommand{\fecha}[1]{\def\Fecha{#1}}
\newcommand{\departamento}[1]{\def\Departamento{#1}}
\newcommand{\carrera}[1]{\def\Carrera{#1}}
% Definir comando para mostrar solo el año actual
\newcommand{\fechaActual}{\number\year}

\newcommand{\caratulaTapa}{
	\begin{titlepage}
		\begin{center}
			{\fontsize{18}{20}\selectfont UNIVERSIDAD CATÓLICA BOLIVIANA SAN PABLO SEDE TARIJA}\\[0.5cm]
			{\fontsize{16}{18}\selectfont DEPARTAMENTO DE \MakeUppercase{\Departamento}}
			
			{\fontsize{14}{16}\selectfont CARRERA DE \MakeUppercase{\Carrera}\\}
			\begin{figure}[h]
				\centering
				\includegraphics[height=7cm]{ucbLOGO}
			\end{figure}
			{\fontsize{16}{18}\selectfont \MakeUppercase{\Titulo} }
			\vspace{1cm}
			
			{\fontsize{14}{16}\selectfont POSTULANTE: \MakeUppercase{\Autor}\\}
		\end{center}
		\vspace{0.2cm}
		{\fontsize{12}{14}\selectfont
			Trabajo de proyecto de grado \textcolor{red}{(lo que corresponda según Reglamento de modalidades de graduación)} presentado en consideración de la Universidad Católica Boliviana "San Pablo, como requisito para optar el Grado Académico de Licenciatura en \Carrera 
		}
	
		{\centering \fontsize{14}{16}\selectfont TARIJA-BOLIVIA\\\Fecha\\}
		
		
		\thispagestyle{empty} % Evita numeración en la portada
	\end{titlepage}
}

% Comando para generar la portada (sin numeración de página)
\newcommand{\caratulaContenido}{
	\begin{titlepage}
		\begin{center}
			{\fontsize{18}{20}\selectfont UNIVERSIDAD CATÓLICA BOLIVIANA SAN PABLO SEDE TARIJA}\\[0.5cm]
			{\fontsize{16}{18}\selectfont DEPARTAMENTO DE \MakeUppercase{\Departamento}}
			
			{\fontsize{14}{16}\selectfont CARRERA DE \MakeUppercase{\Carrera}\\}
			\begin{figure}[h]
				\centering
				\includegraphics[height=7cm]{ucbLOGO}
			\end{figure}
			{\fontsize{16}{18}\selectfont \MakeUppercase{\Titulo} }
			\vspace{1cm}
			
			{\fontsize{14}{16}\selectfont POSTULANTE: \MakeUppercase{\Autor}\\TUTOR: \MakeUppercase{\Tutor}}
		\end{center}
		\vspace{0.2cm}
		{\fontsize{12}{14}\selectfont
			Trabajo de proyecto de grado \textcolor{red}{(lo que corresponda según Reglamento de modalidades de graduación)} presentado en consideración de la Universidad Católica Boliviana "San Pablo, como requisito para optar el Grado Académico de Licenciatura en \Carrera 
		}
	
		{\centering \fontsize{14}{16}\selectfont TARIJA-BOLIVIA\\\Fecha\\}
		
		
		\thispagestyle{empty} % Evita numeración en la portada
	\end{titlepage}
}

% Iniciar numeración desde la introducción
\newcommand{\iniciarNumeracion}{
    \setcounter{page}{1} % Reiniciar la numeración de páginas
    \pagestyle{fancy} % Usar el estilo "fancy"
    \fancyhf{} % Limpiar encabezados y pies de página
    \fancyfoot[R]{\thepage} % Colocar el número de página a la derecha
    \renewcommand{\headrulewidth}{0pt} % Eliminar la línea del encabezado
    \renewcommand{\footrulewidth}{0pt} % Eliminar la línea del pie de página
    % Redefinir el estilo 'plain' para las páginas de inicio de capítulo
    \fancypagestyle{plain}{%
      \fancyhf{}%
      \fancyfoot[R]{\thepage}%
      \renewcommand{\headrulewidth}{0pt}%
      \renewcommand{\footrulewidth}{0pt}%
    }
}
\newcommand{\configurarIndices}{
    % Líder de puntos para el índice general (capítulos, secciones y subsecciones)
    \renewcommand{\cftchapleader}{\cftdotfill{\cftdotsep}}
    \renewcommand{\cftsecleader}{\cftdotfill{\cftdotsep}}
    \renewcommand{\cftsubsecleader}{\cftdotfill{\cftdotsep}}
    
    % Líder de puntos para la lista de figuras
    \renewcommand{\cftfigleader}{\cftdotfill{\cftdotsep}}
    
    % Líder de puntos para la lista de tablas
    \renewcommand{\cfttableader}{\cftdotfill{\cftdotsep}}
}
% Definir datos del documento
\departamento{ciencias de la tecnología e innovación}
\titulo{TITULO DE TRABAJO (escrito en un letra mayuscula redactado en un maximo de 3 lineas alineado y justificado tamaño t16)\\}
\autor{Nombre del Autor}
\tutor{Nombre Tutor}
\fecha{\number\year}
\carrera{ingeniería mecatrónica}
% Generar portada automáticamente
\setcounter{secnumdepth}{5}  % Hasta subparagraph numerado

\begin{document}
	\caratulaTapa
	\caratulaContenido
	
	% Índice general
	\configurarIndices
    \tableofcontents
    \thispagestyle{empty}
    \newpage
    \listoffigures
    \thispagestyle{empty}
    \newpage
    \listoftables
    \thispagestyle{empty}
    \newpage
	
	
	\iniciarNumeracion
	% Secciones principales
	\renewcommand{\thesection}{\arabic{section}}
	\section{ANÁLISIS DEL PROBLEMA}
	Aquí inicia el desarrollo del documento.
	\subsection{Introducción}
	Breve presentación del tema del proyecto y su relevancia
	\subsection{Descripción del problema}
	Detalle especifico del problema que se va a abordar, destacando su importancia.
	\subsection{Formulación del problema}
	Planteamiento claro y conciso del problema del proyecto.
	\subsection{Pregunta de investigación (opcional)}
	Pregunta principal que el proyecto busca responder.
	\subsection{Objetivos del proyecto}
	\subsubsection{Objetivo general}
	Meta  del proyecto.
	\subsubsection{Objetivos específicos}
	Metas secundarias que ayudan a alcanzar el objetivo general.
	\subsection{Delimitación del proyecto}
	Definición del alcance y las limitaciones del estudio.
	\subsection{Justificación}
	Explicación de la relevancia y la importancia del proyecto, y como contribuirá al campo tecnológico.
	
	\newpage
	\section{Marco Teórico}
	\subsection{Antecedentes y referencias}
	Información previa y estudios anteriores relevantes al tema del proyecto.
	\subsection{Estado del arte}
	Resumen y análisis de las investigaciones mas recientes y relevantes en el área, identificando tendencias y vacíos en el conocimiento.
	\subsection{Desarrollo de teorías y modelos}
	Descripción de las teorías y modelos que sustentan el proyecto.
	\subsection{Definición de términos básicos}
	Explicación de términos técnicos y conceptos específicos utilizados en el proyecto
	
	\newpage
	\section{Marco tecnológico}
	
	\subsection{Tecnologías utilizadas y tendencias}
	Identificación de tecnologías, sus tendencias y su impacto potencial en el proyecto.
	
	\subsection{Comparación de tecnologías}
	Análisis comparativo de diferentes tecnologías disponibles.
	
	\newpage
	\section{Marco metodológico}
	
	\subsection{Metodología}
	Descripción detallada del enfoque y métodos utilizados para llevar a cabo el proyecto, incluyendo técnicas de recolección de datos y análisis.
	
	\subsection{Plan de trabajo}
	Cronograma y fases del proyecto, especificando actividades y tiempos estimados para cada etapa.
	\newpage
	\section{Marco legal (opcional)}
	
	\newpage
	\section{Ingeniería del proyecto}
	
	\newpage
	\section{Conclusiones y recomendaciones}
	
	\newpage
	% Bibliografía
	\section*{Bibliografía}
	Aquí van las referencias en formato APA 7ma edición.
	
	\newpage
	% Anexos
	\section*{Anexos}
	Aquí se pueden agregar anexos.
	
	\begin{figure}[ht]
		
		% Incrementa el contador de figuras para numerarlas automáticamente
		\refstepcounter{figure}
		% Número de la figura en negrita
		\textbf{Figura \thefigure}\\[0.5em]
		% Título de la figura en cursiva (una línea debajo del número)
		\textit{Título breve pero descriptivo de la imagen}\\[1em]
		\begin{center}
			
			\includegraphics[width=0.8\textwidth]{eje1.png}\\[1em]
		\end{center}
		% Inserción de la imagen
		
		% Leyenda: explicación de símbolos o detalles de la imagen
		
		% Nota: información adicional si fuese necesaria
		\normalsize Nota: Se incluye la nota únicamente cuando es necesaria para aclarar información adicional.
		\addcontentsline{lof}{figure}{Figura \thefigure. \textit{Título breve pero descriptivo}}
		
	\end{figure}
	
\end{document}